% ============================================================================
% Fichier  : partie1/P01_C04_risque_notions_mesure.tex
% Titre    : Le Risque, Notions et Mesure
% Auteur   : MINKA MI NGUIDJOÏ Thierry Emmanuel
% Version  : 1.0 -- novembre 2026
% ============================================================================
% Base : ACSI C03 (registre LiZbiZ R1 à R8 et appétence repris à
% l'identique ; lecture de la gestion des risques 2/3/5 ; RD ≈ 0,07).
% Sources : ISO 31000:2018 ; IAASB ISA 315 révisée (risque d'anomalies
% significatives) ; INTOSAI GUID 5100 (trois risques d'audit) ; Cox (2008)
% sur les limites des matrices de risque ; Hubbard et Seiersen (2016) pour la
% mesure quantitative ; articles A0 et A1 (incertitude épistémique, NEK).
% Échelles de cotation : convention du cours.
% Corrigés : corriges/Corr_C04.tex.
% ============================================================================

\chapter{Le Risque, Notions et Mesure}
\label{chap:risque}
\chaptermark{Le risque}

\epigraph{%
    « On ne maîtrise pas un risque qu'on n'a pas nommé,\\
    ni un risque qu'on ne sait pas comparer aux autres. »%
}{Principe du gestionnaire de risques}

\niveauChapitre{Fondamental}{%
    Maîtrise du Chapitre~\ref{chap:controle-interne} (contrôle interne,
    matrice des risques et des contrôles)%
}

\begin{boxobjectifs}
\begin{itemize}
    \item Définir le risque, la menace, la vulnérabilité, le risque inhérent,
          le risque résiduel, l'appétence et la tolérance.
    \item Distinguer le risque, qui porte sur les objectifs de
          l'organisation, de l'incertitude épistémique, qui porte sur ce que
          l'auditeur sait.
    \item Coter un risque en probabilité et en impact sur une échelle
          commune, et connaître les limites d'une matrice de risque.
    \item Mesurer un risque en valeur monétaire (perte unitaire, fréquence
          annuelle, perte annuelle attendue) et comparer le coût d'un
          contrôle au risque qu'il réduit.
    \item Appliquer le modèle du risque d'audit pour dimensionner les
          travaux, et l'inverser pour en déduire le risque de détection
          admissible.
\end{itemize}
\end{boxobjectifs}

\bigskip

% ============================================================================
\section*{Question directrice : lequel est le plus grave~?}
% ============================================================================

Le comité de direction de LiZbiZ a reçu, en un mois, trois alertes~: un écart
de stock de 186~millions FCFA, des comptes d'agents partis encore actifs,
et une coupure de courant de douze heures au centre de données de Garoua.
Chaque responsable réclame la priorité pour son dossier. M.~ETEME demande à
Mme~NKENG, la contrôleuse interne~: «~Lequel est le plus grave~?~»

Elle ne peut pas répondre~: LiZbiZ n'a ni liste de ses risques, ni échelle
pour les comparer. Ce chapitre donne les notions et les mesures~; le
Chapitre~\ref{chap:gestion-risques} donnera les méthodes d'analyse et le
processus de gestion.

\medskip
\textbf{Plan du chapitre.}
\begin{enumerate}
    \item Les notions~: risque, composantes, inhérent et résiduel,
          appétence (section~\ref{sec:c04-definir}).
    \item Le risque et l'incertitude de l'auditeur
          (section~\ref{sec:c04-incertitude}).
    \item La mesure qualitative et le registre de LiZbiZ
          (section~\ref{sec:c04-qualitatif}).
    \item Les limites des matrices de risque
          (section~\ref{sec:c04-limites}).
    \item La mesure quantitative (section~\ref{sec:c04-quantitatif}).
    \item Les risques de l'audit (section~\ref{sec:c04-audit}).
\end{enumerate}

\separateur

% ============================================================================
\section{Définir le Risque}
\label{sec:c04-definir}
% ============================================================================

\subsection{Effet de l'incertitude sur les objectifs}

\begin{boxdef}
\termecle{Risque}\index{risque}~: effet de l'incertitude sur les objectifs.
L'effet est un écart, positif ou négatif, par rapport à ce qui est attendu.
On décrit un risque par ses sources, les événements possibles, leurs
conséquences et leur vraisemblance~\cite{iso31000_2018}.
\end{boxdef}
{\small\itshape Définition reformulée d'après l'ISO~31000:2018.}

\medskip

Cette définition a deux conséquences pratiques. D'abord, un risque se
rattache toujours à un \emph{objectif}~: « le centre de Garoua tombe en
panne » n'est un risque qu'au regard d'un objectif, par exemple « suivre
les expéditions du Nord en continu ». Ensuite, un risque peut être aussi
une opportunité~; en audit, on s'intéresse surtout aux écarts défavorables.

\subsection{Écrire un risque}

Un risque mal écrit ne peut être ni coté ni traité. « Le risque
informatique » n'est pas un risque~; « la sécurité » non plus. Le cours
adopte une formulation en trois parties.

\begin{boxprocedure}
\textbf{Formuler un risque~: cause, événement, conséquence}

\textit{En raison de} [cause, vulnérabilité], \textit{il est possible que}
[événement], \textit{ce qui entraînerait} [conséquence sur un objectif].

\medskip
\textbf{Exemple (R3 de LiZbiZ).} En raison de l'absence de procédure de
retrait des comptes au départ d'un salarié, il est possible qu'un ancien
agent utilise son compte encore actif, ce qui entraînerait des opérations
frauduleuses non imputables et une perte financière.
\end{boxprocedure}

\subsection{Les composantes d'un risque informatique}

\begin{boxdef}
\termecle{Menace}\index{menace}~: cause potentielle d'un incident (un employé
malveillant, une panne électrique, un rançongiciel).

\medskip
\termecle{Vulnérabilité}\index{vulnérabilité}~: faiblesse qu'une menace peut
exploiter (un compte non désactivé, l'absence de groupe électrogène).

\medskip
\termecle{Vraisemblance}\index{vraisemblance} ou probabilité~: chance que
l'événement se produise sur une période donnée.

\medskip
\termecle{Impact}\index{impact}~: ampleur des conséquences sur les
objectifs, en argent, en durée d'interruption, en sanction ou en image.
\end{boxdef}

Une menace sans vulnérabilité ne produit rien~: un rançongiciel contre un
système dont les sauvegardes sont hors ligne et testées cause une
interruption, pas une perte de données. Une vulnérabilité sans menace non
plus~: un compte actif d'un ancien salarié qui a quitté le pays présente
peu de risque d'usage par lui-même, mais beaucoup d'usage par un autre.

\subsection{Inhérent, résiduel, appétence, tolérance}

\begin{boxdef}
\termecle{Risque inhérent}\index{risque inhérent}~: niveau du risque avant
la prise en compte des contrôles.

\medskip
\termecle{Risque résiduel}\index{risque résiduel}~: niveau du risque qui
subsiste compte tenu des contrôles réellement en place et efficaces.

\medskip
\termecle{Appétence au risque}\index{appétence au risque}~: niveau de risque
que l'organisation accepte de prendre pour atteindre ses objectifs. Elle est
fixée par l'organe de gouvernance, et non par les opérationnels.

\medskip
\termecle{Tolérance}\index{tolérance au risque}~: écart acceptable autour
d'un objectif précis (par exemple, une disponibilité du système de
facturation d'au moins 99,5~\% par mois). La tolérance traduit
l'appétence en seuils mesurables.
\end{boxdef}

\begin{boxalert}
\textbf{Violation typique.} Calculer le risque résiduel en tenant compte
d'un contrôle décrit dans la procédure, mais jamais testé.

\smallskip
\textbf{Règle~:} on ne réduit le risque que du fait des contrôles dont
l'efficacité est établie. Un contrôle non évalué ne réduit rien~; sa cotation
est \NEK{} tant qu'on ne l'a pas testé (Chapitre~\ref{chap:controle-interne}).
\end{boxalert}

% ============================================================================
\section{Risque et Incertitude de l'Auditeur}
\label{sec:c04-incertitude}
% ============================================================================

Le mot « incertitude » apparaît dans la définition du risque comme dans
celle de l'auditologie (Chapitre~\ref{chap:audit-auditologie}). Il ne
désigne pas la même chose.

\begin{tableaucours}{L{3.2cm}L{4.6cm}L{4.4cm}}{%
    \tableauHeader{} &
    \tableauHeader{Risque} &
    \tableauHeader{Incertitude épistémique}
}
Porte sur & Les événements futurs qui affectent les objectifs de
l'organisation & Ce que l'auditeur (ou l'organisation) sait de l'état
actuel du système \\
Nature & En partie aléatoire~: même parfaitement informé, on ne sait pas si
la panne aura lieu & Réductible~: plus de preuves réduisent l'ignorance \\
Se mesure par & Probabilité et impact & Taux d'incertitude $U$ et
entropie $H$ des réponses \\
Exemple chez LiZbiZ & Le centre de Garoua peut tomber en panne & Personne
ne sait si le groupe électrogène a été testé \\
\end{tableaucours}
\vspace{-0.8em}
\begin{center}{\small\itshape Tableau~4.1. Risque et incertitude
épistémique}\end{center}

La distinction entre composantes aléatoire et épistémique de l'incertitude
est l'une des thématiques du programme de
l'auditologie~\cite{minka2026auditologie}. Elle a une conséquence directe~:
une cotation de risque n'est fiable que si l'on sait ce que valent les
informations qui l'alimentent. L'auditeur cote donc aussi \emph{sa
connaissance} des risques.

\begin{boxlecture}{Gestion des risques de LiZbiZ}
AUDIX CONSEIL pose dix questions sur le dispositif de gestion des risques
(existence d'une liste, d'une échelle, d'un responsable, d'un suivi…) à la
direction et aux responsables de site~: 2~\OUI{}, 3~\NON{}, 5~\NEK{}.

\medskip
\textbf{Constatation.} $\mu = 0{,}20$~; $U = 0{,}50$~; $H \approx 1{,}49$~bit.

\medskip
\textbf{Analyse} (confiance élevée). La moitié des réponses déclare ne pas
savoir~: l'organisation ignore largement comment ses risques sont gérés.
Avant de coter quoi que ce soit, il faut établir qui est responsable de
quel risque. C'est une non-conformité inconnue.

\medskip
\hfill\lecture{0{,}20}{0{,}50}{1{,}49}
\end{boxlecture}

% ============================================================================
\section{La Mesure Qualitative}
\label{sec:c04-qualitatif}
% ============================================================================

\subsection{Une échelle commune}

Pour comparer des risques de natures différentes, on les cote sur une même
échelle. Le cours retient une convention à quatre niveaux, courante en
pratique~; chaque organisation fixe la sienne.

\begin{tableaucours}{cL{3.6cm}L{8.0cm}}{%
    \tableauHeader{Niv.} &
    \tableauHeader{Probabilité (sur un an)} &
    \tableauHeader{Impact}
}
1 & Rare & Mineur~: perte inférieure à 5~millions FCFA, interruption de
moins de 4~heures \\
2 & Possible & Modéré~: de 5 à 50~millions FCFA, interruption de moins de
24~heures \\
3 & Probable & Majeur~: de 50 à 500~millions FCFA, interruption de 1 à
3~jours, réserve du commissaire aux comptes \\
4 & Quasi certain & Critique~: plus de 500~millions FCFA, interruption de
plus de 3~jours, sanction d'un régulateur \\
\end{tableaucours}
\vspace{-0.8em}
\begin{center}{\small\itshape Tableau~4.2. Échelles de cotation (convention
du cours)}\end{center}

Le \termecle{score}\index{score de risque} d'un risque est le produit
$\text{probabilité} \times \text{impact}$, de 1 à 16. Le risque résiduel se
calcule de la même façon, en abaissant la probabilité (ou l'impact) du
seul effet des contrôles efficaces.

\begin{boxinfo}
\textbf{Ancrer les niveaux de probabilité.} « Probable » ne dit rien tant
qu'on ne l'a pas chiffré. Le cours propose~: rare, moins d'une fois en dix
ans (moins de 10~\% par an)~; possible, une fois en trois à dix ans (10 à
30~\%)~; probable, une fois en un à trois ans (30 à 70~\%)~; quasi certain,
au moins une fois par an (plus de 70~\%). Sans ces ancrages, deux personnes
donnent au même risque deux cotes différentes.
\end{boxinfo}

\subsection{Le registre de LiZbiZ}

\begin{boxscenario}{Opération LIZBIZ. Huit risques cotés}
Mme~NKENG réunit les responsables et établit le registre suivant. La colonne
« Maîtrise » indique de combien de niveaux les contrôles \emph{testés}
abaissent la probabilité.

\medskip
\begin{tableaucours}{cL{4.6cm}ccccc}{%
    \tableauHeader{Réf.} &
    \tableauHeader{Risque} &
    \tableauHeader{P} & \tableauHeader{I} & \tableauHeader{Inh.} &
    \tableauHeader{Maîtr.} & \tableauHeader{Rés.}
}
R1 & Ajustement frauduleux du stock (cumul réception et ajustement) & 4 & 3
& 12 & 0 & 12 \\
R2 & Paiement de fournisseurs fictifs & 3 & 3 & 9 & 0 & 9 \\
R3 & Usage du compte d'un agent parti & 3 & 3 & 9 & 0 & 9 \\
R4 & Arrêt prolongé du centre de Garoua & 2 & 4 & 8 & 1 & 4 \\
R5 & Perte de données faute de sauvegarde restaurable & 2 & 4 & 8 & 0 & 8 \\
R6 & Fuite de données personnelles des clients & 2 & 3 & 6 & 1 & 3 \\
R7 & Erreurs de facturation & 3 & 2 & 6 & 1 & 4 \\
R8 & Chiffrement des serveurs par un rançongiciel & 2 & 4 & 8 & 0 & 8 \\
\end{tableaucours}

\medskip
\textbf{Appétence fixée par le conseil.} Un score résiduel de 9 ou plus est
intolérable et doit être traité sans délai~; de 5 à 8, il est traité dans
l'année~; de 1 à 4, il est accepté sous surveillance.

\medskip
\textbf{Constatation.} Trois risques sont intolérables (R1, R2, R3), deux sont
à traiter dans l'année (R5, R8), trois sont acceptés (R4, R6, R7).

\medskip
\textbf{Analyse} (confiance moyenne). Les trois risques intolérables
relèvent tous du contrôle interne (séparation des tâches, gestion des
comptes), pas de la technique~: ils confirment le diagnostic du
Chapitre~\ref{chap:controle-interne}. La coupure de Garoua, qui inquiétait
le plus le comité, est en réalité le risque le mieux maîtrisé.
\end{boxscenario}

\subsection{Coter avec plusieurs évaluateurs}

Une cote est un jugement. Pour réduire la part d'arbitraire, AUDIX
CONSEIL fait coter chaque risque séparément par trois personnes, puis
compare. Lorsque les trois cotes de probabilité diffèrent de plus d'un
niveau, la cote n'est pas retenue~: on cherche l'information qui manque.

\begin{boxscenario}{Opération LIZBIZ. Trois évaluateurs pour R8}
\textbf{Faits.} Pour R8 (rançongiciel), M.~KOUNOU cote la probabilité~1
(«~nous avons un antivirus~»), Mme~NKENG~3, le consultant d'AUDIX CONSEIL~2.

\medskip
\textbf{Analyse.} L'écart (de 1 à 3) révèle une incertitude épistémique sur
la vulnérabilité, pas un désaccord d'opinion~: personne ne sait si les
sauvegardes sont isolées du réseau. La cote~2 est retenue provisoirement,
avec la mention~: \emph{cotation à confirmer par le test des sauvegardes}
(Chapitre~\ref{chap:audit-continuite}).
\end{boxscenario}

% ============================================================================
\section{Les Limites des Matrices de Risque}
\label{sec:c04-limites}
% ============================================================================

La matrice probabilité × impact est l'outil le plus répandu. Elle a des
défauts connus, que l'auditeur doit savoir nommer lorsqu'il évalue la
cartographie d'une organisation~\cite{cox2008matrices}.

\begin{tableaucours}{L{3.6cm}L{8.6cm}}{%
    \tableauHeader{Limite} &
    \tableauHeader{Illustration}
}
Résolution faible & Un risque à 6~millions et un risque à 49~millions ont
le même impact~2~; leurs scores ne disent pas leur écart \\
Inversion possible & Deux risques de même score (8) peuvent avoir des
pertes attendues très différentes~: P2 × I4 et P4 × I2 ne représentent pas
la même chose \\
Échelle ordinale & Le produit de deux rangs (2 × 4) n'est pas une
grandeur~: un score de 8 n'est pas « deux fois » un score de 4 \\
Corrélations ignorées & Deux risques notés séparément peuvent se réaliser
ensemble (le rançongiciel et la perte des sauvegardes) \\
Biais des évaluateurs & Optimisme du propriétaire du risque, ancrage sur
le dernier incident \\
\end{tableaucours}
\vspace{-0.8em}
\begin{center}{\small\itshape Tableau~4.3. Cinq limites des matrices de
risque}\end{center}

\begin{boxalert}
\textbf{Violation typique.} Additionner les scores des risques d'un domaine
pour en déduire un « risque global » (« le domaine informatique totalise
45~points »).

\smallskip
\textbf{Règle~:} les scores ordinaux ne s'additionnent pas. Pour agréger,
on compte les risques par zone de l'appétence, ou l'on passe à une mesure
en valeur (section~\ref{sec:c04-quantitatif}).
\end{boxalert}

% ============================================================================
\section{La Mesure Quantitative}
\label{sec:c04-quantitatif}
% ============================================================================

\subsection{Perte unitaire, fréquence, perte annuelle attendue}

Lorsque l'organisation dispose de données (coût horaire d'une interruption,
fréquence des incidents passés), le risque se mesure en valeur. Le modèle
élémentaire, largement enseigné, repose sur trois grandeurs.

\begin{boxdef}
\termecle{Perte unitaire}\index{perte unitaire} (\textit{single loss
expectancy}, SLE)~: perte causée par une occurrence de l'événement,
$\text{SLE} = \text{valeur de l'actif} \times \text{facteur d'exposition}$.

\medskip
\termecle{Fréquence annuelle}\index{fréquence annuelle}
(\textit{annualized rate of occurrence}, ARO)~: nombre moyen d'occurrences
par an (0,25 signifie une fois en quatre ans).

\medskip
\termecle{Perte annuelle attendue}\index{perte annuelle attendue}
(\textit{annualized loss expectancy}, ALE)~:
$\text{ALE} = \text{SLE} \times \text{ARO}$.
\end{boxdef}

\begin{boxscenario}{Opération LIZBIZ. Le centre de Garoua en chiffres}
\textbf{Données} (DAF et DSI)~: le centre de Garoua traite les expéditions
du Nord, qui génèrent une marge de 6~millions FCFA par heure d'activité. Un
arrêt prolongé type dure 4~jours (96~heures) faute de pièces de rechange~;
on retient un facteur d'exposition de 1 (toute la marge est perdue pendant
l'arrêt). Sans groupe électrogène testé, la DSI estime un tel arrêt
possible une fois en trois ans environ (ARO~$\approx 0{,}3$).

\medskip
\textbf{Calcul.}
\begin{align*}
\text{SLE} &= 96 \times 6 = 576~\text{millions FCFA}, \\
\text{ALE}_{\text{avant}} &= 576 \times 0{,}3 \approx 173~\text{millions FCFA par an}.
\end{align*}
Le groupe électrogène, une fois testé chaque mois, ramène l'ARO à 0,1~:
$\text{ALE}_{\text{après}} \approx 58$~millions FCFA par an.

\medskip
\textbf{Analyse.} La réduction de la perte attendue, environ 115~millions
par an, justifie un contrôle (groupe électrogène, contrat de maintenance,
test mensuel) coûtant jusqu'à ce montant. La cote qualitative d'impact~4
est confirmée (plus de 500~millions par occurrence).
\end{boxscenario}

\subsection{Le coût d'un contrôle}

Le calcul précédent donne une règle de décision simple, qu'on appelle
parfois valeur de la mesure de sécurité~:
\[
    \text{valeur du contrôle} = \text{ALE}_{\text{avant}}
    - \text{ALE}_{\text{après}} - \text{coût annuel du contrôle}.
\]
Un contrôle de valeur négative coûte plus qu'il ne rapporte~: l'organisation
peut décider d'accepter le risque, par une décision écrite.

\subsection{Raisonner en intervalles}

Les valeurs ponctuelles (6~millions par heure, 0,3 par an) cachent leur
imprécision. La pratique recommandée est d'estimer chaque grandeur par un
intervalle de confiance à 90~\%, dont on est sûr à 90~\% qu'il contient la
vraie valeur~\cite{hubbard2016cyber}. Le
Chapitre~\ref{chap:gestion-risques} combinera ces intervalles par
simulation de Monte-Carlo (méthode FAIR).

\begin{boxinfo}
\textbf{Un exercice de calibration.} Les évaluateurs non entraînés donnent
des intervalles trop étroits~: ils sont trop sûrs d'eux. On les calibre en
leur posant des questions dont on connaît la réponse (la population de
Douala, la date de création de la CEMAC) et en leur montrant combien de
leurs intervalles contiennent la bonne réponse. Un évaluateur calibré à
90~\% voit environ neuf intervalles justes sur dix.
\end{boxinfo}

\subsection{Qualitatif ou quantitatif~?}

\begin{tableaucours}{L{2.8cm}L{4.6cm}L{4.8cm}}{%
    \tableauHeader{} &
    \tableauHeader{Qualitatif} &
    \tableauHeader{Quantitatif}
}
Données requises & Jugements d'experts & Coûts, fréquences, durées \\
Résultat & Rangs, zones de couleur & Montants, distributions de pertes \\
Agrégation & Impossible au sens strict & Possible (somme des pertes
attendues, simulation) \\
Usage & Premier tri, dialogue avec les métiers & Arbitrage budgétaire,
décision sur un contrôle coûteux \\
\end{tableaucours}
\vspace{-0.8em}
\begin{center}{\small\itshape Tableau~4.4. Mesure qualitative et mesure
quantitative}\end{center}

Les deux ne s'opposent pas~: on trie avec la matrice, on chiffre les
risques qui demandent une décision d'investissement.

% ============================================================================
\section{Les Risques de l'Audit}
\label{sec:c04-audit}
% ============================================================================

\subsection{Le modèle du risque d'audit}

L'auditeur court lui-même un risque~: conclure qu'une situation est
satisfaisante alors qu'elle ne l'est pas. La profession décompose ce
\termecle{risque d'audit}\index{risque d'audit} en trois éléments, que le
GUID~5100 de l'INTOSAI reprend pour l'audit des SI~\cite{intosai2019guid5100}.

\begin{boxdef}
\termecle{Risque inhérent}~(en audit)~: probabilité qu'une anomalie
significative existe, du fait de la nature même de l'activité ou du système.

\medskip
\termecle{Risque de contrôle}\index{risque de contrôle}~: probabilité que le
contrôle interne ne prévienne pas ou ne détecte pas cette anomalie.

\medskip
\termecle{Risque de détection}\index{risque de détection}~: probabilité que
les travaux de l'auditeur ne la détectent pas.
\end{boxdef}

Le modèle classique relie les trois~:
\[
    \text{RA} = \text{RI} \times \text{RC} \times \text{RD}.
\]
L'auditeur ne choisit ni le risque inhérent ni le risque de contrôle~: il
les évalue. Leur produit est ce que la norme ISA~315 révisée appelle le
risque d'anomalies significatives~\cite{iaasb2019isa315}. L'auditeur
choisit en revanche le risque de détection, en dimensionnant ses tests,
pour que le risque d'audit reste au niveau qu'il accepte. On inverse donc
le modèle~:
\[
    \text{RD} = \frac{\text{RA}}{\text{RI} \times \text{RC}}.
\]

\begin{boxscenario}{Opération LIZBIZ. Dimensionner les tests sur les stocks}
\textbf{Évaluation.} Risque inhérent élevé (stock de forte valeur, écart
déjà constaté), estimé à 0,9~; risque de contrôle élevé (contrôles de stock
inefficaces, Chapitre~\ref{chap:controle-interne}), estimé à 0,8. Risque
d'audit accepté~: 5~\%.

\medskip
\textbf{Calcul.} $\text{RD} = 0{,}05 / (0{,}9 \times 0{,}8) \approx 0{,}07$.

\medskip
\textbf{Analyse.} Le risque de détection doit être très faible~: AUDIX
CONSEIL ne peut pas se contenter de tester les contrôles. Il doit réaliser
des tests directs et étendus sur les données de stock
(Chapitre~\ref{chap:audit-applications}).
\end{boxscenario}

\subsection{Du risque de détection à l'étendue des tests}

Le Tableau~4.5 traduit le niveau du risque de détection admissible en
nature des travaux. C'est une grille de raisonnement, pas un barème.

\begin{tableaucours}{L{2.8cm}L{4.4cm}L{5.0cm}}{%
    \tableauHeader{RD admissible} &
    \tableauHeader{Nature des travaux} &
    \tableauHeader{Étendue}
}
Élevé (plus de 0,5) & Tests des contrôles, procédures analytiques &
Échantillons réduits, revue globale \\
Moyen (0,2 à 0,5) & Tests des contrôles et tests directs ciblés &
Échantillons standards (Chapitre~\ref{chap:preuve}) \\
Faible (moins de 0,2) & Tests directs étendus, analyse de toute la
population & Extraction complète, requêtes sur toutes les opérations \\
\end{tableaucours}
\vspace{-0.8em}
\begin{center}{\small\itshape Tableau~4.5. Du risque de détection aux
travaux}\end{center}

\begin{boiteRemarque}{Des probabilités vraiment mesurables~?}
Rarement. Les valeurs 0,9 ou 0,8 traduisent un jugement, pas une mesure.
Le modèle sert surtout à raisonner~: plus le risque inhérent et le risque de
contrôle sont élevés, plus les tests doivent être étendus. L'auditologie
propose d'aller plus loin en mesurant ce que l'auditeur sait réellement,
par $U$ et $H$, plutôt que de l'estimer (Chapitre~\ref{chap:conclure}).
\end{boiteRemarque}

\subsection{Un audit fondé sur les risques}

La cartographie des risques de l'organisation oriente le plan d'audit~:
on audite d'abord ce qui est le plus risqué et le moins bien connu. Deux
dimensions guident donc le choix des zones d'audit~: le score résiduel
(section~\ref{sec:c04-qualitatif}) et l'incertitude sur ce score
(section~\ref{sec:c04-incertitude}). Un risque coté~8 avec une forte
incertitude, comme R8, peut passer avant un risque coté~9 bien documenté. Le
Chapitre~\ref{chap:planifier} montrera comment passer de la cartographie
au plan d'audit.

\begin{boxmetier}{Gestionnaire de risques}
\textbf{Tenir une cotation défendable}
\begin{enumerate}
    \item Faire approuver l'échelle et ses ancrages chiffrés par le conseil.
    \item Écrire chaque risque en trois parties (cause, événement,
          conséquence).
    \item Faire coter par plusieurs personnes et investiguer les écarts.
    \item Ne réduire le résiduel que sur la base de contrôles testés.
    \item Chiffrer en valeur les risques qui appellent un investissement.
\end{enumerate}
\end{boxmetier}

\begin{boxmetier}{Auditeur SI}
\textbf{Auditer la cotation, pas seulement la cartographie.} L'auditeur
demande, pour chaque risque rouge, la justification de la probabilité et de
l'impact, et le rapport de test des contrôles qui ont réduit le résiduel.
Une cotation sans justification se cote \NEK{}.
\end{boxmetier}

\begin{boxinfo}
\textbf{Ce que dit la recherche.} L'article A0~\cite{minka2026auditologie}
sépare l'incertitude aléatoire (ce qui arrivera) et l'incertitude
épistémique (ce que l'on sait). L'article A1~\cite{minka2026nek} en tire
la mesure $U$ de la seconde. Ce chapitre applique la distinction au
registre~: une cote est un fait sur le risque \emph{et} un fait sur la
connaissance de celui qui cote.
\end{boxinfo}

\separateur

\begin{boxresume}
\textbf{Ce qu'il faut retenir de ce chapitre~:}
\begin{itemize}
    \item \textbf{Le risque est l'effet de l'incertitude sur les objectifs}~:
          on l'écrit en trois parties (cause, événement, conséquence).
    \item \textbf{Inhérent, résiduel, appétence, tolérance}~: le résiduel ne
          tient compte que des contrôles testés~; l'appétence est fixée par
          la gouvernance et traduite en tolérances.
    \item \textbf{Risque et incertitude épistémique} ne se confondent pas.
    \item \textbf{La mesure qualitative} trie~; elle a des limites (faible
          résolution, échelle ordinale, corrélations ignorées)~; ses scores
          ne s'additionnent pas.
    \item \textbf{La mesure quantitative} (SLE, ARO, ALE) chiffre le risque
          et la valeur d'un contrôle~; on estime par intervalles.
    \item \textbf{Le risque d'audit} = inhérent × contrôle × détection~;
          l'auditeur en déduit le risque de détection admissible et
          l'étendue de ses tests.
\end{itemize}
\end{boxresume}

\begin{boxinfo}
\textbf{Points de contrôle de l'élève.} Je sais~: écrire un risque en trois
parties~; coter sur l'échelle du cours~; nommer trois limites des
matrices~; calculer SLE, ARO, ALE~; inverser le modèle du risque d'audit.
\end{boxinfo}

\bigskip

\begin{boxexo}[title={\ding{45}\quad Exercice~4.1. Menace, vulnérabilité, impact \niveaubase}]
Pour chacun des risques R3, R5 et R8 du registre de LiZbiZ, identifiez la
menace, la vulnérabilité et l'objectif menacé, puis justifiez la cotation
d'impact au regard du Tableau~4.2.
\end{boxexo}

\begin{boxexo}[title={\ding{45}\quad Exercice~4.2. Écrire un risque \niveaubase}]
Réécrivez en trois parties (cause, événement, conséquence) les risques
suivants, formulés par les responsables de LiZbiZ~: (a) «~Le réseau~»~;
(b) «~Les mots de passe~»~; (c) «~Le prestataire de Garoua~».
\end{boxexo}

\begin{boxexo}[title={\ding{45}\quad Exercice~4.3. Après les contrôles \niveaumoyen}]
LiZbiZ met en place les contrôles proposés au Chapitre~\ref{chap:controle-interne}.
Les tests montrent qu'ils abaissent de deux niveaux la probabilité de R1 et
de R3, et d'un niveau celle de R2.
\begin{enumerate}[(a)]
    \item Recalculez les scores résiduels de R1, R2 et R3.
    \item Quels risques restent à traiter dans l'année~?
    \item Pourquoi faut-il attendre les tests avant de mettre à jour le
          registre~?
\end{enumerate}
\end{boxexo}

\begin{boxexo}[title={\ding{45}\quad Exercice~4.4. Deux risques de score 8 \niveaumoyen}]
R4 (P2, I4) et un risque fictif R9 (P4, I2) ont le même score inhérent.
\begin{enumerate}[(a)]
    \item En prenant le milieu de chaque tranche de l'échelle (probabilité
          et perte), estimez la perte annuelle attendue de chacun.
    \item Qu'en concluez-vous sur le score~?
\end{enumerate}
\end{boxexo}

\begin{boxexo}[title={\ding{45}\quad Exercice~4.5. La valeur d'un contrôle \niveaumoyen}]
Reprenez le centre de Garoua. Le groupe électrogène, son contrat de
maintenance et le test mensuel coûtent 40~millions FCFA par an.
\begin{enumerate}[(a)]
    \item Calculez la valeur du contrôle.
    \item À partir de quelle marge horaire le contrôle cesse-t-il d'être
          rentable, toutes choses égales par ailleurs~?
    \item Une solution de secours en nuage ramènerait la durée d'un arrêt
          à 8~heures, pour 70~millions par an, sans changer l'ARO initial
          de 0,3. Est-elle préférable~?
\end{enumerate}
\end{boxexo}

\begin{boxexo}[title={\ding{45}\quad Exercice~4.6. Trois évaluateurs \niveaumoyen}]
Pour R5, les trois cotes de probabilité sont 1, 2 et 4. Que faites-vous~?
Quelle information cherchez-vous, et auprès de qui~? Comment le registre
doit-il présenter R5 en attendant~?
\end{boxexo}

\begin{boxexo}[title={\ding{45}\quad Exercice~4.7. Inverser le modèle \niveaumoyen}]
Pour l'audit des paiements de LiZbiZ, AUDIX CONSEIL estime RI~=~0,7 et,
après test des contrôles, RC~=~0,3. Le risque d'audit accepté est de
5~\%.
\begin{enumerate}[(a)]
    \item Calculez le risque de détection admissible.
    \item Quelle nature et quelle étendue de travaux le Tableau~4.5
          suggère-t-il~?
    \item Comparez avec les stocks. Pourquoi les deux zones ne reçoivent-
          elles pas le même effort~?
\end{enumerate}
\end{boxexo}

\begin{boxexo}[title={\ding{45}\quad Exercice~4.8. Additionner des scores \niveaumoyen}]
Le rapport annuel de la DSI de la BANQUE SANAGA indique~: «~Le risque
informatique global est passé de 52 à 47~points, soit une amélioration de
10~\%.~» Rédigez la remarque de l'auditeur.
\end{boxexo}

\begin{boxexo}[title={\ding{45}\quad Exercice~4.9. Une échelle commune pour la BANQUE SANAGA \niveauexpert}]
\begin{enumerate}[(a)]
    \item Proposez une échelle d'impact commune aux risques bancaires et
          informatiques, exprimée en FCFA et en durée d'interruption.
    \item Cotez sur cette échelle le risque « échec d'une migration du
          système central ».
    \item Évaluez le risque inhérent et le risque de contrôle d'un audit de
          la gestion des changements de la banque, et déduisez-en l'étendue
          des tests.
    \item Expliquez pourquoi une réponse \NEK{} sur l'existence d'un plan de
          retour arrière doit être signalée au conseil.
\end{enumerate}
\end{boxexo}

\begin{boxexo}[title={\ding{45}\quad Exercice~4.10. Intervalles et perte attendue \niveauexpert}]
Pour R8 (rançongiciel), la DSI estime la perte unitaire entre 100 et
900~millions FCFA (intervalle à 90~\%) et la fréquence entre 0,05 et 0,4
par an.
\begin{enumerate}[(a)]
    \item Calculez une perte annuelle attendue «~basse~» et «~haute~» en
          combinant les bornes. Pourquoi cet encadrement est-il trop
          large~?
    \item Expliquez pourquoi une simulation donnerait une réponse plus
          utile (Chapitre~\ref{chap:gestion-risques}).
    \item Quelle information réduirait le plus l'incertitude~?
\end{enumerate}
\end{boxexo}

\bigskip

\begin{boxinfo}
\textbf{Pour aller plus loin.}
\begin{itemize}
    \item ISO~31000:2018~\cite{iso31000_2018}~: le vocabulaire et les
          principes, en quelques pages.
    \item L.~A.~Cox, \textit{What's Wrong with Risk
          Matrices?}~\cite{cox2008matrices}~: la démonstration des
          limites des matrices.
    \item D.~Hubbard et R.~Seiersen, \textit{How to Measure Anything in
          Cybersecurity Risk}~\cite{hubbard2016cyber}~: la mesure par
          intervalles et la calibration.
    \item IAASB, ISA~315 révisée~\cite{iaasb2019isa315}~: le risque
          d'anomalies significatives dans l'audit des comptes.
\end{itemize}
\end{boxinfo}

\begin{boxinfo}
\textbf{Transition.} Les notions et les mesures sont en place. Le
Chapitre~\ref{chap:gestion-risques} les assemble dans un processus de
gestion des risques et présente les méthodes d'analyse utilisées en
pratique (EBIOS Risk Manager, MEHARI, OCTAVE, FAIR).
\end{boxinfo}
